\bookchapter{第三十三章：未来展望——スマホの次のステップ}

\chapterintro{スマホは一夜で消えないが、その“独占的地位”は次世代デバイスに分流される可能性がある。この章では三つのことを語る：次世代デバイスはどんな姿か、スマホがいかに“隠身”して新しいアイデンティティの基盤になるか、そして私たちが前もって考えるべき倫理的議題。}

\booksection{一、AR/VRと空間コンピューティング：スマホ以外の“第二のデスクトップ”}

次世代デバイスの競争はすでに始まっており、スマホの時間を最も分流しそうなのはAR/VRと空間コンピューティングデバイスだ：

\begin{itemize}
  \item \textbf{Vision Pro後継}：Apple Vision Proの発表は"空間コンピューティング"概念の正式な実装を示したが、初代製品の価格（29999元から）と重量は依然ハードル。2026-2027年により軽くて安い普及版Visionが出ると予想され、空間コンピューティングが"開発者のおもちゃ"から"初期マス層"へと移行する。その時、iPhoneはVisionグラスの演算パートナーとデータハブになる可能性がある。
  \item \textbf{折りたたみスマートフォン}：目の前に"第二のデスクトップ"を開き、デジタル情報が空気のように至る所に存在するように。MetaとRay-Banのコラボスマートグラスはすでに数百万個売れ、国内のRokid、XREAL、INMOなどもコンシューマー向けARグラスを推出。2025-2026年のARグラスはリアルタイム翻訳、ナビHUD、AIアシスタント対話ができるが、光導波ディスプレイとバッテリーの制限があり、"スマホ代替"までまだ3-5年ある。
  \item \textbf{スマートリングとウェアラブル}：Samsung Galaxy Ring、Oura Ringなどのスマートリングが流行し始め、スマホの"肌身につけたセンサー"として健康モニタリングを補完。スマホの役割は"万能端末"から"ハブ+複数のライト化アクセサリ"の組み合わせへと移行している。
\end{itemize}

\booksection{二、ブレイン·マシン·インターフェース：思考がインタラクションの媒介に}

思考自体がインタラクションの媒介になるが、この道のりは想像より長い：

\begin{itemize}
  \item \textbf{Neuralink最新の進展}：MuskのNeuralinkは2024年に初の人体埋め込みを完了し、最初の被験者Noland Arbaughが思考でコンピュータのマウスを操作してチェスを指した。現在、埋め込み型BCI（ブレイン·マシン·インターフェース）は主に麻痺患者の補助シーンに向けられ、彼らが再びデジタル世界を"操作"できるようにすることが目標だ。
  \item \textbf{非侵襲型BCI}：複数の企業が非侵襲型ヘッドバンド（Kernel、OpenBCIなど）を開発し、手術なしで脳波信号を読み取れるが、精度は埋め込み型に遥かに劣る。"思考でタイピング"までにはまだかなりの距離がある。
  \item \textbf{コンシューマー向けの展望}：ブレイン·マシン·インターフェースがコンシューマー向けスマホインタラクションに到達するまであと10-20年だが、短期的にはアクセシビリティ補助や集中力モニタリングなどのシーンで先に使われるかもしれない。倫理とプライバシーの問題（"私の脳波は誰のプライバシーか"）が技術の成熟に先んじて議論の焦点になる。
\end{itemize}

\booksection{三、スマホの“隠身”と新しいアイデンティティの基盤}

スマホのもう一つの未来の進化方向は“隠身”かもしれない——至る所に存在する“黒子”となり、あなたのデジタルアイデンティティの唯一の基盤になる：

\begin{itemize}
  \item \textbf{二、スマートフォンの「不可視化」と新たなアイデンティティの礎}：瞬時に"あなたはあなた"の認証を完了し、あらゆるサービスのログインで手動認証コード入力が不要に。
  \item \textbf{クラウド二要素認証}：新規スマホ起動時にデジタル資産継承方案を事前設定し、"身後の事"を出荷時デフォルトに。
  \item \textbf{デジタル遺産が初期設定に}：あなたのクラウド上の個人データ金庫こそが中心で、ハードウェアはアクセスの入口にすぎない。機種変更、デバイス変更しても、アイデンティティとデータは常にあなたについてくる。
\end{itemize}

\booksection{四、倫理的思考：私たちはまだスマホの主人か？}

：あなたのクラウド上の個人データ金庫こそが中心。

\begin{itemize}
  \item \textbf{AI倫理}：オンデバイスAIが自分よりも好みを理解する時、あなたの「選択」は自主的な意志か、誘導された結果か？AI Agentが代わりに注文やメッセージ返信をする時、責任と境界はどこにありますか？
  \item \textbf{データ主権}：あなたの健康データ、脳波データ、生体特徴データは、あなたに帰属するのか、デバイスメーカーに帰属するのか、クラウドサービスプロバイダーに帰属するのか？「個人情報保護法」とEUのGDPRは出発点にすぎず、「データポータビリティ権」と「忘れられる権利」の真の実装にはまだ長い道のりがあります。
  \item \textbf{デジタル·ディバイド}：AIスマホ、折りたたみ画面、衛星通信はすべて"使える"と"使えない"の格差を広げている。デジタルインフラへの公平なアクセスは、20年前の"ネットが使えるか"と同じように重要な新たな社会課題になる。
  \item 三、倫理的考察：私たちは、まだスマートフォンの主人なのか？
  \item 一つの究極の問いが常に存在します：私たちはテクノロジーによって力を与えられた「神」なのか、それともアルゴリズムに餌付けされた「水槽の脳」なのか？
\end{itemize}

\booksection{五、繋がりを保ち、独立も保つ——あなたのデジタルライフ哲学}

技術がどう進化しようと、以下のデジタルライフ哲学は10年後も適用できる：

\begin{enumerate}
  \item 数分おきに無意識に画面を点灯させているとき、それはあなたの習慣ですか、それとも設計された「依存行動」ですか？
  \item 四、つながりを保ちながら、独立も保つ——あなたのデジタルライフ哲学
  \item リソースの分配者であり、プッシュ通知の受け手ではない。
  \item 自動化をあなたの忠実な「コパイロット」にする。
  \item あなたの堅牢な「デジタル後方拠点」を構築する。
\end{enumerate}

\begin{tcbpractice}{定期的に「デジタル断捨離」と「ソーシャル断食」を実行する。}
あなたとスマホ（そして将来のARグラスやAIアシスタント）のために簡単な“利用規約”を書き、関係を定義する。例えば：

\begin{itemize}
  \item 試してみよう：あなたの「個人デジタルデバイス利用サービス規約」を作成する
  \item あなたとあなたのスマートフォンのために簡単な「サービス規約」を書き、関係を定義しましょう。例えば：
  \item 第1条：私の注意の時間は私のものです。あなたは私が仕事/家族と過ごす時に自ら「サイレントモード」に入るべきです。
\end{itemize}

\end{tcbpractice}

\begin{tcbsummary}{本章のまとめ}
\begin{itemize}
  \item Vision ProとARグラスはスマホ以外の新たなコンピューティング入口だが、短期的にはスマホを代替できない
  \item ブレイン·マシン·インターフェースはコンシューマー向けまであと10-20年、倫理問題が技術の成熟に先んじる
  \item スマホは"隠身"方向に進化し、デバイス中心からアイデンティティ中心へ
  \item AI倫理、データ主権、デジタル·ディバイドは今後10年の核心議題
  \item 繋がりを保ち独立も保つことは、デジタルライフの不変の哲学
\end{itemize}

\end{tcbsummary}
